\section{The private set}
\label{app:heldout}

Every task of the track has a public answer, and every task of its stream a public solver, so the statement of what a learner may use, which comes with the code, forbids the lookup of Appendix~\ref{app:bench} but cannot make it impossible. A new seed changes the instances and not the rules. The private set holds new rules that only we have seen, on which we rerun a learner's code to verify a result. We keep it as ARC-AGI keeps its private evaluation set \citep{chollet2025arcagi2}: the code comes to the tasks, and the tasks never go to a submitter's machine, which could keep whatever it was shown. We name no rule here and show no instance, since anything shown of a rule publishes it. The construction below took the first track's set from 98 candidates to the 50 kept; the second and third tracks' sets, built by the same method, follow it.

\paragraph{The rules.} Six agents on Claude Opus 5.5 wrote the 98 candidates, one agent for each family of rules: objects, geometry, lines and paths, patterns and sequences, panels and logic, and physics and topology, with 16 or 17 rules each. A rule is a module with a generator, which draws an instance from a seed, and a solver, which computes the output from the input alone. The benchmark serves the rules as one more task source, with a fresh pair for every test input and every demonstration, so any learner that runs on the public track runs on the private set unchanged. Its held-out rules are generated as well, and a request there returns a fresh pair up to the cap of eight, where the public held-out test serves each task's own two to seven pairs. The agents were asked for one idea that a person could infer from three examples, with every color role, count or direction that the answer depends on visible in the grid, one answer for every instance, varied sizes, positions and colors, and no public task or plain ARC trope, such as gravity or symmetry completion, without a twist that changes the rule.

\paragraph{Filters and screens.} A candidate passes the mechanical filters when, over 300 draws, its generator and solver always finish, every grid is valid, the solver reproduces every generated output, its inputs and outputs vary, and the output is not always the input under one fixed rotation or reflection; all 98 passed, and the thresholds are in the set's code. Two screens then look for a rule that already exists. Re-ARC's 400 verifiers, the solvers of the ARC-AGI-1 training tasks, reproduced none of the candidates, and each candidate's own solver, given every pair of all 1{,}920 public ARC-AGI-1 and ARC-AGI-2 tasks, dropped six candidates, each of which matched one to three public tasks. No candidate contained another.

\paragraph{Calibration.} We measure how often a model solves a rule at first sight. An item is one fresh instance with three demonstrations, drawn as the demo-cache floor first meets the task, in a fixed prompt, and the model gives one answer, which must match exactly, with no feedback and no second attempt. Claude Opus 5.5 at max effort answered four instances of each of the 92 remaining candidates and solved 360 of the 368. GPT-6 Luna answered eight of each at low, medium and high effort. Both answered the ARC-AGI-1 track's 50 public tasks in the same way: fresh instances of its 25 stream tasks from the Re-ARC generators, as many as for a rule, and each test pair of its 25 held-out evaluation tasks with up to three of the task's own pairs. A candidate that Opus solved fewer than three times in four was set aside, which two were, and one that Luna solved every time at every effort would have been too easy, which none was. The calibration cost about 60 US dollars on Opus and 6 on Luna. These rates are not the floors' first-exposure rates of Appendix~\ref{sec:models}, since a floor chooses its own demonstration requests and often answers before it asks for any.

\paragraph{Selection.} From the 90 eligible candidates, we removed one at a time, each time the one whose removal brought the rest nearest the public tasks' first-sight rates over the four presets of Claude Opus 5.5 and GPT-6 Luna, keeping the families near equal shares, until 50 were left, 8 or 9 per family, and split them into 25 for the stream and 25 held out, each family dealt to both halves (Table~\ref{tab:heldout}). The 92 calibrated candidates were easier for GPT-6 Luna than the public tasks, by 3.7 to 5.6 points, and the selection removes the difference. The other 40 eligible candidates are a reserve, for a kept rule that has to be retired. The set is frozen, since a run on it is comparable with another only while its rules stay as they are. It serves one user, since every further user brings 25 rules of their own, and we will add rules to verify learners in the collective regimes.

\paragraph{Two models the selection never saw.} After the set was frozen, GLM-5.3 Flash and Qwen3.5-9B at medium effort answered eight instances of each of the 92 candidates and the ARC-AGI-1 track's public tasks in the same way, as a check on models that the selection did not match. GLM-5.3 Flash solves 26.5\% of the kept rules and 26.8\% of the public tasks, though its halves part in opposite directions, the streamed rules 6.5 points under the public stream and the held-out rules 6.0 points over the public held-out tasks. Qwen3.5-9B solves 5.5\% of the kept rules and 8.5\% of the public tasks, its streamed rules about as often as the public stream and its held-out rules 6.5 points less often than the public held-out tasks. The public held-out rates rest on 26 answers, one to each test pair, of which GLM-5.3 Flash solved 6 and Qwen3.5-9B 3, so one answer moves them by up to four points. More than nine in ten of Qwen3.5-9B's replies, on the rules and on the public tasks alike, reached the cap of 32{,}000 tokens before an answer, so its rates mostly count replies cut off. GLM-5.3 Flash's reached it on 22\% of its replies on the rules and 25\% on the public tasks, almost none of them solved, so its rates also count some replies cut off, about equally on both sides. The two cost about 9 US dollars together.

\begin{table}[htbp]
\centering
\small\setlength{\tabcolsep}{5pt}
\begin{tabular}{lrrrrrrr}
\toprule
 & \multicolumn{2}{c}{stream} & \multicolumn{2}{c}{held out} & \multicolumn{2}{c}{all} & \\
\cmidrule(lr){2-3}\cmidrule(lr){4-5}\cmidrule(lr){6-7}
model, effort & private & public & private & public & private & public & calibrated \\
\midrule
GPT-6 Luna, low & 31.5 & 31.5 & 21.0 & 20.0 & 26.3 & 25.8 & 31.4 \\
GPT-6 Luna, medium & 54.0 & 54.0 & 56.0 & 56.0 & 55.0 & 55.0 & 58.7 \\
GPT-6 Luna, high & 59.0 & 59.0 & 66.5 & 66.0 & 62.8 & 62.5 & 67.7 \\
Claude Opus 5.5, max & 97.0 & 97.0 & 100.0 & 100.0 & 98.5 & 98.5 & 97.8 \\
GLM-5.3 Flash, medium & 23.0 & 29.5 & 30.0 & 24.0 & 26.5 & 26.8 & 29.2 \\
Qwen3.5-9B, medium & 5.5 & 5.0 & 5.5 & 12.0 & 5.5 & 8.5 & 4.9 \\
\bottomrule
\end{tabular}

\caption{First-sight solve rates, in percent, on the private set and on the ARC-AGI-1 track's public tasks, measured the same way: three demonstrations, one answer and an exact match. Private columns are the 25 streamed rules, the 25 held-out rules and the 50 kept; public columns are the track's 25 stream tasks, its 25 held-out ARC-AGI-1 evaluation tasks and all 50. Calibrated is all 92 candidates that passed the filters and screens. Each rate is a mean over rules or tasks, of eight instances per rule or stream task for GPT-6 Luna, GLM-5.3 Flash and Qwen3.5-9B and four for Claude Opus 5.5. The selection matched the first four rows; the last two were calibrated after it.}
\label{tab:heldout}
\end{table}

\paragraph{The second track's set.} Seven agents on Claude Opus 5.5 wrote 116 candidate rules for the season, one agent for each of seven families new to the track: legends and keys, context, compositions, programs and procedures, relations and graphs, assembly and shape matching, and counting and reading out. Each was asked for ARC-AGI-2's character: symbols whose meaning the grid sets, rules that compose or apply by context, and grids up to 30 by 30. All 116 passed the filters and both screens, against the 1{,}920 public tasks and fresh draws of the track's ARC-GEN generators, and none contains another or one of the first track's rules. Claude Opus 5.5 at max effort solved 450 of its 464 answers, and the two rules it solved only twice in four were set aside. The selection matched Kimi K3.1 Pro at medium effort, the track's model, with 24 answers per rule, and kept 50 rules, 7 per family and 8 of counting, 25 for the stream and 25 held out, with 64 in reserve. Kimi solves 32\% of the kept rules and of the public tasks, the stream 45 against 45 and the held-out rules 18 against 18, and Opus 95 against 91. Two models the selection never saw agree: Kimi-K3 at medium effort solves 38\% against 40 and GLM-5.3 Flash 12 against 12, though half or more of GLM-5.3 Flash's replies reach the cap. These readings are of OpenRouter's routing on 7 November 2026, which served Kimi from many providers; pinned to Moonshot AI's own endpoint, which served all our Kimi runs, Kimi solves 30\% of the kept rules against 27 of the public tasks, the stream 46 against 41 and the held-out rules 15 against 14. On eight fresh instances of each kept rule, which the selection never read, Kimi still solves 32\%, but its halves part less, the stream 38 and the held-out rules 26, so some of the spread the selection matched was chance, and the held-out rules may read easier than the public held-out tasks. The calibration cost about 660 US dollars.

\paragraph{The third track's set.} ARC-AGI-3 keeps no hidden games that anyone outside ARC Prize can use, so agents on Claude Opus 5.5 wrote 28 candidate games for the season on ARC Prize's engine, each with sprite types and mechanics of its own, a first level, a generator and a solver, and each designed away from the 25 public games. Several are versions of well-known puzzles, and none is an ARC-AGI-3 game. All 28 pass the public track's checks at 50 seeds, draw the same levels in another process, and draw without failing at 500 of the seeds a run derives. A first-sight reading plays fresh levels as Appendix~\ref{app:tracks} does, the model wiped before each level and free to ask for three demonstrations. Claude Opus 5.5 at max effort solved all 112 levels it played, and all 100 of the public games'. Kimi K3.1 Pro at medium effort, the track's model, first solved fewer levels of the candidates than of the public games, 19 against 28\% on the first 22, and the eight of which it solved none, puzzles that need planning ahead, were made easier with their rules unchanged, with smaller boards, fewer pieces or shorter first levels, before the selection. The selection kept 25 games, 15 for the stream and 10 held out, with 3 in reserve, on which Kimi solves 28\% of the levels against 28 on the public games, the stream 31 against 31 and the held-out games 24 against 24, on OpenRouter's routing of 7 November. That match is selected on the readings that measure it. Pinned to Moonshot AI's own endpoint, which served all our Kimi runs and serves every verified run, the same levels read the kept games somewhat harder, 27\% against 33, the stream 29 against 35 and the held-out games 24 against 29, a gap of about one standard error over games. GLM-5.3 Flash at medium effort, which the selection never saw, solves 43\% against 40 on the routing, the held-out games 50 against 42, within the noise of 40 levels. The kept games are mostly played by clicking, 15 by clicks, 8 by keys and 2 by both, where the public games are 8, 4 and 13. A learner no better than its reset reads a held-out share saved within about $[-0.14, 0.12]$ on the 10 held-out games, against $[-0.11, 0.10]$ on the public track's. No person has played the games. The calibration cost about 580 US dollars.

\paragraph{The verified run.} A submitter opens an issue on the benchmark's public repository with a public commit, the learner, the published models it uses and a budget for their calls, which it pays ahead. We run the learner, its reset and each of its models bare on the private set at a seed drawn for the run and not disclosed, inside the sandbox of Appendix~\ref{app:harness} on a machine of its own that has no public network and holds neither the rules nor the key, and the learner reaches nothing but the tracking proxy, which makes every model call on a key of ours, each model served by one endpoint for the whole run, capped at that budget. The issue then receives each test's share saved with its 95\% interval, the share learned, the budgets and shares solved of the learner and the reset, and what the model calls cost, with the flags below. Only these numbers come back. The records, and anything the learner kept, stay with us, since they show the rules.

\paragraph{Checks against gaming.} A submission names every model its learner uses, each being a published model that must start as released. Each share is read beside its budgets (Table~\ref{tab:tests-absolute}), and a test on which the reset plays below any of the learner's models played bare at no greater cost, or below the cheapest of them at any cost, is flagged, so that a deliberately weakened reset cannot inflate the learner's relative gain, and so is one on which the learner plays no better than such a model, so that switching to a stronger model cannot read as learning, and one on which the learner sends a larger share of its calls than its reset to any of its models but the cheapest, so that switching to a stronger model for only some calls cannot either.

\paragraph{What remains open.} No person has tried the rules. Claude Opus 5.5 at max effort, solving every kept rule at least three times in four from three demonstrations, stands in for the human testers who solved each ARC-AGI-2 task \citep{chollet2025arcagi2}. The same model wrote the rules, and since it solves 98.5\% of both the kept rules and the public tasks, at the ceiling, the calibration cannot show whether rules written by a model are easier for that model. An Opus learner is read against itself reset on the same instances, which shares any such advantage. The match of Table~\ref{tab:heldout} is at first sight and says nothing of how much a rule can be learned on its return, so the private set's learnable budget need not equal the public stream's, and a verified run measures it with the reset. We read a share learned on the private set below the one on the public stream, beyond both intervals, as a sign of tuning to the public stream. A learner chooses its own scores, so the numbers a verified run publishes carry a little of what it was served. They are rounded, a submitter's runs in a season are capped, and the set can be retired for new rules, but only rules drawn fresh for every run would remove this. Learners that train their weights are verified on machines we rent, since a submitter's own machine would keep the instances it trained on.
